\documentclass[twocolumn]{aastex631}
\begin{document}
\title{A residual reionization ionizing-photon-budget shortfall for star-forming galaxies at z\~{}11 (optimistic systematic corner)}
\author{NebulaMind Lab (autonomous pipeline)}
\affiliation{NebulaMind Lab, \url{https://lab.nebulamind.net}}
\begin{abstract}
Reionization ionizing-photon-budget reconciliation at z\~{}11: star-forming galaxies require f\_esc=0.276 (+0.238/-0.127) to close the budget (Madau-Dickinson SFRD, log xi\_ion=25.5+/-0.15, clumping C=2-5, JWST-SFRD tail), versus indirect-proxy-inferred f\_esc=0.080 (+0.147/-0.051) from LzLCS O32/beta calibrations. Median delta(required-inferred)=+0.174 dex-frac (16-84\%: +0.008 to +0.414); 85\% of the systematic MC shows a shortfall. Conclusion: a genuine SHORTFALL remains, and the sign holds under both O32 and beta calibrations. The result is bounded by the xi\_ion x clumping x proxy-calibration systematic, not by statistics. Generated autonomously from public data (jwst) via the NebulaMind Lab runner: a bounded, reproducible, descriptive study.
\end{abstract}
\section{Introduction}
Recent studies have highlighted a potential crisis in the photon budget during reionization, suggesting that current estimates of ionizing photons from star-forming galaxies may not be sufficient to drive this process [Muñoz2024]. This discrepancy has sparked interest in reconciling the ionizing-photon-budget using literature-anchored calculations. Previous works have explored various aspects of reionization, including the role of absorption-dominated scenarios [Davies2021] and excursion set models [Park2022], as well as assessments of galaxy ionizing photon budgets at lower redshifts [Duncan2015]. The analytic approach to cosmic reionization presented in Madau (2017) provides a useful framework for understanding the underlying physics.
\section{Data and method}
To address this issue, we employed a systematics reconciliation over published literature values. Specifically, we adopted the cosmic star formation rate density (SFRD) from Madau \& Dickinson (2014), along with previously published calibrations for the ionizing photon production efficiency (xi\_ion) and escape fraction (f\_esc) proxy based on O32/beta ratios [LzLCS: Chisholm+22, Flury+22; Simmonds+24]. Our method focused on calculating the ionizing-photon-budget at z\~{}11 using these parameters.
\section{Result}
Our analysis reveals that star-forming galaxies require an escape fraction of f\_esc = 0.276 (+0.238/-0.127) to reconcile the reionization photon budget, given the Madau-Dickinson SFRD, log xi\_ion = 25.5 +/- 0.15, and clumping factor C=2-5. In contrast, indirect proxy-inferred values from LzLCS O32/beta calibrations yield f\_esc = 0.080 (+0.147/-0.051). The median difference between required and inferred escape fractions is +0.174 dex-frac (16-84\%: +0.008 to +0.414), with 85\% of systematic Monte Carlo simulations showing a shortfall.
\begin{figure}\centering\includegraphics[width=\columnwidth]{result.png}\caption{A residual reionization ionizing-photon-budget shortfall for star-forming galaxies at z\~{}11 (optimistic systematic corner)}\end{figure}
\section{Caveats}
It is essential to acknowledge the limitations of our approach, which relies on automated, single-selection, uncalibrated measurements from published literature. The accuracy of our result depends heavily on the assumptions and calibrations used in previous studies, such as the choice of xi\_ion and f\_esc proxy. Additionally, uncertainties in clumping factor C and SFRD introduce further variability into our calculations. These limitations emphasize the need for more direct observations and refined models to better constrain the reionization photon budget.
\section*{References}
\footnotesize
[Muoz2024] Muñoz et al. (2024). Reionization after JWST: a photon budget crisis?. bibcode 2024MNRAS.535L..37M \par [Davies2021] Davies et al. (2021). The Predicament of Absorption-dominated Reionization: Increased Demands on Ionizing Sources. bibcode 2021ApJ...918L..35D \par [Park2022] Park et al. (2022). Calibrating excursion set reionization models to approximately conserve ionizing photons. bibcode 2022MNRAS.517..192P \par [Duncan2015] Duncan et al. (2015). Powering reionization: assessing the galaxy ionizing photon budget at z < 10. bibcode 2015MNRAS.451.2030D \par [Madau2017] Madau (2017). Cosmic Reionization after Planck and before JWST: An Analytic Approach. bibcode 2017ApJ...851...50M

\end{document}
